% ===================== CHƯƠNG 1: TỔNG QUAN =========================
\chapter{TỔNG QUAN VỀ ĐỀ TÀI}
\label{ch:tongquan}

Chương này phân tích và đánh giá các vấn đề liên quan trực tiếp đến đề tài: bối cảnh bài toán dự báo bỏ học MOOC, các họ phương pháp đã được đề xuất, nội dung và phương pháp nghiên cứu của từng họ, và cuối cùng là tổng hợp khoảng trống nghiên cứu định hướng cho phần đóng góp của đồ án.

\section{Bối cảnh và phát biểu vấn đề}
\label{sec:boicanh}

\subsection{Hiện tượng bỏ học trong MOOC}

MOOC ra đời với kỳ vọng dân chủ hóa giáo dục đại học, nhưng thực tế vận hành lại bộc lộ một vấn đề nan giải về khả năng duy trì người học. Theo nhận định của de Freitas và cộng sự, mặc dù MOOC mở rộng đáng kể khả năng tiếp cận, mô hình này đồng thời bộc lộ tỷ lệ duy trì và tỷ lệ hoàn thành rất thấp \cite{deFreitas2015}. Khảo sát có hệ thống của Deng và cộng sự trên các trải nghiệm người học MOOC giai đoạn 2008--2021 xác nhận rằng các nền tảng lớn chỉ ghi nhận 5--15\% người đăng ký hoàn thành khóa, và phần lớn rời bỏ trong vài tuần đầu \cite{Deng2022}.

Hiện tượng rời bỏ sớm này có hai hệ quả quan trọng đối với việc thiết kế hệ thống dự báo. Một mặt, nó nhấn mạnh tầm quan trọng của việc \emph{cảnh báo sớm}: nếu phần lớn người học rời bỏ trong vài tuần đầu, thì một hệ thống chỉ có thể dự báo chính xác sau khi đã thu thập đủ dữ liệu của nhiều tuần sẽ trở nên vô dụng về mặt can thiệp. Mặt khác, nó tạo ra một bài toán dữ liệu mất cân bằng nghiêm trọng, khi nhóm người học có nguy cơ cao hoặc đã rời khóa chiếm tỷ trọng áp đảo.

\subsection{Dự báo bỏ học trong bối cảnh phân tích học tập}

Bài toán dự báo bỏ học không tồn tại đơn lẻ mà nằm trong lĩnh vực rộng hơn là \textit{phân tích học tập} (\textit{learning analytics}) -- việc thu thập, đo lường, phân tích và báo cáo dữ liệu về người học và bối cảnh học tập nhằm hiểu và tối ưu hóa quá trình học. Trong khung này, dự báo bỏ học là một bài toán \emph{dự báo} (predictive analytics) đóng vai trò đầu vào cho một bài toán \emph{quy định} (prescriptive analytics) là thiết kế can thiệp. Mối quan hệ này giải thích vì sao việc chỉ tối ưu một chỉ số phân loại đơn lẻ là chưa đủ: giá trị cuối cùng của một mô hình dự báo nằm ở chỗ nó cung cấp tín hiệu đúng lúc và đủ chi tiết để hành động. Đây là lý do đồ án nhấn mạnh đồng thời ba khía cạnh -- thời điểm cảnh báo, độ chi tiết của nhãn, và chất lượng phát hiện lớp can thiệp -- thay vì chỉ chạy theo một con số Accuracy đơn lẻ.

\subsection{Vì sao cần phân biệt nhiều mức rủi ro}

Cách tiếp cận truyền thống đặt bài toán dưới dạng phân loại nhị phân: người học hoặc hoàn thành, hoặc không hoàn thành. Tuy nhiên, cách gộp nhóm này che giấu một sự phân biệt có ý nghĩa can thiệp cốt lõi. Trong nhóm ``không hoàn thành'', cần phân biệt rõ giữa người học \emph{đã thực sự rời khóa} (không còn khả năng can thiệp) với người học \emph{vẫn đang hoạt động nhưng có nguy cơ cao} (vẫn còn cơ hội can thiệp). Hai nhóm này đòi hỏi chiến lược hoàn toàn khác nhau: với nhóm đầu, can thiệp gần như vô ích; với nhóm sau, một nhắc nhở kịp thời có thể làm thay đổi quỹ đạo học tập. Việc gộp chung hai nhóm dưới một nhãn ``không hoàn thành'' do đó làm giảm đáng kể giá trị thực tiễn của mô hình dự báo. Đây chính là động cơ cho việc đặt bài toán dưới dạng nhãn rủi ro có thứ tự nhiều mức.

Hơn nữa, một nhãn năm mức còn cho phép phân biệt giữa các mức hoàn thành tốt (A, B, C), điều có giá trị cho các bài toán liên quan như xếp hạng, cấp chứng chỉ hay phân luồng nội dung. Mặc dù đồ án tập trung vào lớp can thiệp D, việc giữ đủ năm mức bảo toàn thông tin cho các ứng dụng mở rộng về sau và phản ánh trung thực hơn phổ kết quả học tập thực tế.

\section{Tiêu chí khảo sát tài liệu}
\label{sec:tieuchi}

Để bảo đảm tính khách quan và tính đại diện, tài liệu khảo sát trong đồ án được chọn theo bốn tiêu chí:

\begin{enumerate}[label=(\roman*),leftmargin=1.2cm,topsep=3pt]
  \item liên quan trực tiếp đến dự báo bỏ học hoặc cảnh báo người học có nguy cơ trong môi trường MOOC hoặc học trực tuyến quy mô lớn;
  \item có mô tả đủ rõ về dữ liệu, mô hình hoặc thước đo để cho phép đối sánh khoa học;
  \item đại diện cho các họ phương pháp khác nhau;
  \item ưu tiên các nghiên cứu gốc thay vì bản tóm lược thứ cấp.
\end{enumerate}

Theo các tiêu chí trên, phần khảo sát bao gồm năm nghiên cứu tổng quan có hệ thống nền tảng \cite{Dalipi2018,Prenkaj2020,Aljohani2025,Deng2022,deFreitas2015}, hai bộ dữ liệu chuẩn \cite{Yu2021,Kuzilek2017}, và mười sáu công trình thực nghiệm tiêu biểu từ 2014 đến 2025. Các nghiên cứu tổng quan đóng vai trò định hình bức tranh tổng thể: Dalipi và cộng sự khảo sát các kỹ thuật học máy giai đoạn 2014--2018 \cite{Dalipi2018}; Prenkaj và cộng sự mở rộng đến 2020 và chỉ ra rằng phần lớn công trình tập trung vào trích xuất đặc trưng thủ công từ nhật ký nhấp chuột, so sánh các mô hình học máy truyền thống và đánh giá theo nhãn nhị phân \cite{Prenkaj2020}; gần đây hơn, Aljohani và cộng sự chỉ ra ba xu hướng hiện tại của lĩnh vực gồm chuyển dịch sang học sâu cho dữ liệu chuỗi, ứng dụng meta-heuristic và tối ưu hóa siêu tham số, và tích hợp khả năng giải thích (XAI) \cite{Aljohani2025}.

\subsection{Sự tiến hóa của lĩnh vực qua ba giai đoạn}

Nhìn tổng thể trên trục thời gian, lĩnh vực dự báo bỏ học MOOC đã trải qua ba giai đoạn phát triển có thể nhận diện rõ. \textit{Giai đoạn nền móng (2014--2017)} đặc trưng bởi việc áp dụng các mô hình học máy truyền thống trên đặc trưng thủ công, với đóng góp chính là chứng minh rằng dropout có thể dự báo được từ tín hiệu hành vi \cite{Kloft2014,Taylor2014,Fei2015,Whitehill2017}. \textit{Giai đoạn học sâu (2018--2022)} chứng kiến sự chuyển dịch sang các mô hình chuỗi và kiến trúc lai, cùng với các nghiên cứu phản biện cho thấy học sâu không phải luôn vượt trội \cite{Xing2019,Mubarak2021,Fu2021,Basnet2022}. \textit{Giai đoạn tinh vi hóa (2023--2025)} mở rộng bài toán theo nhiều hướng đồng thời: nhãn đa mức, mô hình đồ thị, tối ưu hóa meta-heuristic và khả năng giải thích \cite{Martins2023,Junejo2025,Aouarib2025,Nguyen2025,Benabbes2025}. Sự tiến hóa này cho thấy một xu hướng rõ: bài toán ngày càng được đặt gần với nhu cầu thực tiễn (phân biệt mức rủi ro, cảnh báo sớm, giải thích được), và đây chính là quỹ đạo mà đồ án tiếp nối và đẩy xa thêm bằng cách hợp nhất nhiều khía cạnh vào một khung benchmark thống nhất.

Các công trình thực nghiệm được phân tích thành hai họ phương pháp lớn dưới đây, sau đó tổng hợp khoảng trống ở tiểu mục cuối.

\subsection{Hai bộ dữ liệu chuẩn của lĩnh vực}

Trước khi đi vào các họ phương pháp, cần giới thiệu hai bộ dữ liệu được sử dụng phổ biến nhất, vì lựa chọn dữ liệu ảnh hưởng trực tiếp đến khả năng đối sánh giữa các công trình.

\textit{MOOCCubeX} \cite{Yu2021} là một kho dữ liệu quy mô lớn hướng tri thức, được thiết kế cho học thích ứng trong MOOC. Bộ dữ liệu bao gồm khoảng 3,3 triệu người học, hơn 4.000 khóa học và hàng trăm triệu sự kiện hành vi, cùng với một đồ thị tri thức liên kết các khái niệm. Quy mô và độ chi tiết của MOOCCubeX khiến nó trở thành lựa chọn lý tưởng cho các nghiên cứu cần cohort lớn và tín hiệu hành vi phong phú -- đây là lý do đồ án chọn MOOCCubeX làm dữ liệu chính.

\textit{OULAD} (Open University Learning Analytics Dataset) \cite{Kuzilek2017} là bộ dữ liệu phân tích học tập của Open University, bao gồm thông tin về sinh viên, khóa học, đánh giá và nhật ký tương tác với môi trường học ảo. OULAD nhỏ hơn MOOCCubeX về quy mô nhưng có ưu điểm là cấu trúc rõ ràng và được sử dụng rộng rãi, do đó thường được dùng để kiểm chứng tính tổng quát. Tuy nhiên, OULAD không phát hành trực tiếp dưới dạng nhãn năm mức A/B/C/D/E, nên việc mở rộng đồ án sang OULAD đòi hỏi một quy tắc ánh xạ nhãn cần được kiểm chứng cẩn thận -- một điểm sẽ được bàn trong phần giới hạn.

Sự khác biệt giữa hai bộ dữ liệu này cũng giải thích một phần vì sao các công trình hiện có khó đối sánh trực tiếp: chúng được đánh giá trên các dữ liệu khác nhau, với các quy tắc nhãn và giao thức chia tập khác nhau. Đây chính là một trong những động cơ thúc đẩy việc xây dựng một khung benchmark thống nhất.

\section{Họ phương pháp 1: Học máy truyền thống và tối ưu hóa}
\label{sec:hoML}

\subsection{Các mô hình nền móng}

Hướng tiếp cận sớm nhất, đại diện bởi Kloft và cộng sự năm 2014, sử dụng SVM tuyến tính trên các đặc trưng thủ công trích từ nhật ký nhấp chuột theo tuần \cite{Kloft2014}. Cùng năm, Taylor và cộng sự đề xuất hồi quy logistic với bộ đặc trưng tương tự để dự báo hiện tượng ``ngừng học'' (\textit{stopout}) \cite{Taylor2014}. Mặc dù đơn giản, các công trình này đặt nền móng quan trọng cho việc xem dropout là một bài toán có thể dự báo từ tín hiệu hành vi -- một tiền đề mà toàn bộ lĩnh vực về sau dựa vào.

Một phát hiện đáng chú ý đến từ Alamri và cộng sự năm 2019: ngay cả khi chỉ dùng hai đặc trưng dễ thu thập trong tuần đầu, các mô hình cây quyết định tăng cường như RF, AdaBoost, GradientBoosting hay XGBoost vẫn đạt độ chính xác 82--94\% trên 21 đợt mở của 5 khóa FutureLearn \cite{Alamri2019}. Kết quả này hàm ý rằng một mô hình cảnh báo rất sớm có thể rất nhẹ về mặt tính toán -- một quan sát có ý nghĩa lớn đối với việc thiết kế hệ thống thực tế. Vitiello và cộng sự bổ sung bằng chứng cho thấy các đặc trưng phản ánh mức độ tham gia tích cực mang tín hiệu dự báo mạnh \cite{Vitiello2018}.

\subsection{Tích hợp khả năng giải thích}

Một biến thể gần đây của họ này tích hợp các công cụ giải thích nhằm tăng độ tin cậy và khả năng thuyết phục người dùng. Nagrecha và cộng sự xây dựng một pipeline có thể diễn giải nhằm thuyết phục giảng viên về kết quả dự báo \cite{Nagrecha2017}; Dass và cộng sự kết hợp RF với SHAP, đạt độ chính xác khoảng 87,5\% trên khóa học tự định nhịp (\textit{self-paced}) \cite{Dass2021}; Benabbes và cộng sự đề xuất một pipeline mô-đun hóa kết hợp XGBoost với SHAP và LIME \cite{Benabbes2025}. Tuy nhiên, hướng tích hợp giải thích này nằm ngoài phạm vi đóng góp chính của đồ án và chỉ được nhắc đến như một hướng đối sánh.

\subsection{Tối ưu hóa lựa chọn đặc trưng và siêu tham số}

Song song với việc cải tiến mô hình, một nhánh quan trọng tập trung vào tối ưu hóa không gian tìm kiếm. Psathas và cộng sự năm 2021 sử dụng Grid Search và Bayesian Optimization để tinh chỉnh nhiều bộ phân loại trên dữ liệu MOOC smart-city \cite{Psathas2021}. Đại diện hiện đại nhất của nhánh này là Aouarib và cộng sự năm 2025: kết hợp giải thuật di truyền (GA) để đồng thời chọn đặc trưng và siêu tham số cho XGBoost trên bộ dữ liệu OULAD, đánh giá tại tuần 5/15/25/35 và báo cáo cải thiện nhất quán so với LR, MLP, DT, SVM, RF và LSTM \cite{Aouarib2025}. Hướng tối ưu hóa cho thấy rằng học máy truyền thống vẫn còn nhiều dư địa cải thiện thông qua thiết kế tốt hơn của không gian tìm kiếm, thay vì phụ thuộc hoàn toàn vào năng lực biểu diễn của mô hình sâu.

\subsection{Phân tích chi tiết phương pháp luận của họ thứ nhất}

Để hiểu rõ vì sao các mô hình họ này có một trần hiệu năng nhất định, cần phân tích sâu hơn về phương pháp luận chung của chúng. Điểm chung xuyên suốt là quy trình hai bước: trước hết, dữ liệu chuỗi hành vi thô được \emph{tổng hợp} thành một vector đặc trưng cố định chiều (ví dụ tổng số lần xem video trong toàn bộ cửa sổ, trung bình điểm bài tập, độ lệch chuẩn của thời gian học); sau đó, một bộ phân loại dạng bảng được huấn luyện trên vector này. Bước tổng hợp, dù làm cho bài toán trở nên đơn giản và nhanh, lại là một ``nút thắt thông tin'': nó làm phẳng chiều thời gian, xóa bỏ thông tin về \emph{trình tự} và \emph{động học} của hành vi. Hai người học có cùng tổng số giờ học nhưng một người tăng dần và một người giảm dần sẽ có cùng vector đặc trưng tổng hợp, dù quỹ đạo rủi ro của họ trái ngược nhau.

Đây chính là lý do sâu xa giải thích vì sao, như sẽ thấy trong kết quả thực nghiệm, các mô hình họ này gần như không cải thiện khi cửa sổ quan sát kéo dài: thông tin bổ sung của các tuần sau bị nén lại trong cùng một phép tổng hợp và không được khai thác như một chuỗi. Quan sát này không phủ nhận giá trị của họ phương pháp truyền thống -- chúng vẫn vượt trội về tốc độ và khả năng giải thích -- mà chỉ định vị rõ ranh giới năng lực của chúng, từ đó biện minh cho sự cần thiết của các mô hình chuỗi thời gian khi mục tiêu là khai thác động học theo tuần.

\section{Họ phương pháp 2: Học sâu và tiếp cận dựa trên đồ thị}
\label{sec:hoDL}

\subsection{Mô hình chuỗi thời gian}

Họ phương pháp thứ hai chuyển trọng tâm từ đặc trưng thủ công sang mô hình hóa trực tiếp chuỗi hành vi theo tuần. Fei và Yeung năm 2015 so sánh ba kiến trúc chuỗi -- IOHMM, RNN và LSTM -- trên dữ liệu Coursera và edX, trong đó LSTM đạt hiệu năng cao nhất \cite{Fei2015}. Đây là một trong những công trình đầu tiên khẳng định giá trị của mô hình tuần tự trong bài toán này. Whitehill và cộng sự đào sâu phân tích các đặc trưng thời gian và xác nhận rằng tín hiệu hành vi theo tuần là một tín hiệu dự báo mạnh \cite{Whitehill2017}. Xing và Du năm 2019 đặt trọng tâm vào việc ước lượng xác suất bỏ học cá nhân hóa bằng mạng nơ-ron sâu, hướng tới mục tiêu can thiệp cá nhân hóa \cite{Xing2019}.

Các kiến trúc lai sau đó kết hợp nhiều thành phần để tăng năng lực biểu diễn. Mubarak và cộng sự đề xuất mô hình CNN + LSTM trên ma trận tương quan 30 ngày, đồng thời tích hợp cơ chế xử lý mất cân bằng vào hàm mất mát \cite{Mubarak2021}; Fu và cộng sự đề xuất CLSA (CNN + LSTM + Attention) với hiệu năng cao trên bộ KDD Cup 2015 \cite{Fu2021}.

\subsection{Một kết luận cảnh báo: học sâu không phải lúc nào cũng vượt trội}

Một đóng góp có tính phản biện quan trọng đến từ Basnet và cộng sự năm 2022. Thông qua đối sánh có hệ thống giữa học máy truyền thống và học sâu trên cùng dữ liệu, nhóm tác giả chỉ ra rằng học máy trong nhiều trường hợp có thể ngang bằng hoặc thậm chí vượt học sâu \cite{Basnet2022}. Kết luận này nhấn mạnh rằng việc lựa chọn mô hình cần dựa trên bằng chứng thực nghiệm trên cùng một giao thức, thay vì giả định mặc nhiên rằng mô hình phức tạp hơn sẽ tốt hơn -- một quan điểm mà đồ án này kế thừa và củng cố. Yu và cộng sự mở rộng hướng học sâu bằng kỹ thuật meta-transfer learning giữa các MOOC khác chủ đề, cho thấy khả năng chuyển giao tri thức giữa các khóa học \cite{Yu2022}.

\subsection{Mở rộng sang nhãn đa mức và mô hình đồ thị}

Nhánh phát triển mới và quan trọng nhất đối với đồ án này là việc mở rộng bài toán từ nhãn nhị phân sang nhãn đa mức, đôi khi kết hợp với mô hình đồ thị để khai thác quan hệ giữa người học. Martins và cộng sự năm 2023 dự báo ba mức graduate/enrolled/dropout trong bối cảnh giáo dục đại học \cite{Martins2023}; Junejo và cộng sự năm 2025 đề xuất 1D-CNN cho nhãn bốn mức trên OULAD \cite{Junejo2025}; Vlachopoulos và cộng sự tổng quan có hệ thống 70 bài về dự báo đa lớp dưới điều kiện mất cân bằng, đồng thời xác nhận rằng các nghiên cứu dự báo đa lớp ``vẫn còn tương đối hiếm'' \cite{Vlachopoulos2025}.

Trong số các công trình sử dụng nhãn đa mức, Nguyen và cộng sự là một trong số hiếm hoi đặt bài toán ở năm mức A/B/C/D/E trên MOOCCubeX, kết hợp GCN adversarial với GRU, và báo cáo kết quả mạnh ($\FoneD = 0{,}92$, độ chính xác $= 0{,}95$) \cite{Nguyen2025}. Đây là công trình gần nhất về mặt thiết kế bài toán với đồ án hiện tại. Tuy vậy, công trình này còn ba hạn chế chưa được giải quyết: (i) chỉ đánh giá tại một cửa sổ cố định (4 tuần); (ii) chiến lược điền giá trị thiếu (GCN adversarial) được đặt cố định thay vì benchmark song song với các lựa chọn thay thế; và (iii) chi phí dựng đồ thị có thể đạt $O(n^2)$, gây khó khăn khi mở rộng quy mô. Schwerter và cộng sự bổ sung thêm một cảnh báo về phương pháp luận: các mô hình phụ thuộc vào quan hệ toàn cục giữa người học có xu hướng kém tổng quát khi chuyển giữa các khóa học hoặc cơ sở đào tạo \cite{Schwerter2026}. Chính những hạn chế này tạo ra không gian mà EarlyRisk-MOOC hướng tới lấp đầy.

\subsection{Phân tích chi tiết phương pháp luận của họ thứ hai}

Họ phương pháp học sâu chia sẻ một triết lý đối lập với họ thứ nhất: thay vì tổng hợp chuỗi rồi phân loại, chúng giữ nguyên cấu trúc chuỗi và để mô hình tự học cách trích xuất đặc trưng thời gian. Tuy nhiên, trong nội bộ họ này cũng có sự phân hóa đáng kể về mức độ phù hợp giữa kiến trúc và dữ liệu.

Các mô hình hồi quy thuần (RNN, LSTM, BiLSTM, GRU) xử lý chuỗi một cách tự nhiên và là lựa chọn phù hợp nhất cho dữ liệu hành vi tuần. Các mô hình lai CNN + LSTM \cite{Mubarak2021,Fu2021} bổ sung khả năng trích xuất đặc trưng cục bộ qua tích chập trước khi đưa vào lớp hồi quy; cách tiếp cận này hiệu quả khi dữ liệu có cấu trúc không gian (ví dụ ma trận tương quan ngày-hành động), nhưng độ phức tạp bổ sung không phải lúc nào cũng được biện minh. Ở thái cực khác, một MLP làm phẳng chuỗi \cite{Xing2019} -- dù vẫn là mô hình ``sâu'' -- thực chất đánh mất cấu trúc tuần tự ngay từ đầu vào, do đó không tận dụng được lợi thế của học sâu trên dữ liệu chuỗi. Sự phân hóa này củng cố một nguyên tắc thiết kế quan trọng: hiệu năng không đến từ ``độ sâu'' của mô hình mà từ \emph{sự phù hợp} giữa kiến trúc và cấu trúc dữ liệu.

Về nhánh đồ thị, sức mạnh của các mô hình GCN nằm ở khả năng khai thác quan hệ giữa người học -- ví dụ những người học cùng khóa hoặc có hành vi tương tự có thể cung cấp thông tin bổ sung cho nhau. Tuy nhiên, sức mạnh này đi kèm ba cái giá: chi phí dựng đồ thị $O(n^2)$, sự phụ thuộc vào chất lượng của đồ thị được dựng, và -- quan trọng nhất theo Schwerter và cộng sự \cite{Schwerter2026} -- xu hướng kém tổng quát khi quan hệ toàn cục giữa người học thay đổi giữa các khóa hoặc cơ sở. Đồ án này đặt giả thuyết rằng trong bối cảnh cửa sổ thời gian sớm, một chiến lược điền giá trị thiếu cục bộ và nhẹ (KNN) có thể đạt hiệu năng tương đương hoặc tốt hơn các phương pháp đồ thị toàn cục mà không phải trả các cái giá trên -- giả thuyết sẽ được kiểm chứng trong Chương~\ref{ch:ketqua}.

\section{Phân tích hiện trạng và khoảng trống nghiên cứu}
\label{sec:khoangtrong}

Bảng~\ref{tab:khaosat} tổng hợp 19 công trình tiêu biểu theo một bộ tiêu chí thống nhất gồm: nhãn sử dụng, phương pháp, dữ liệu, hiệu năng báo cáo, cửa sổ quan sát, ưu điểm và hạn chế.

\begin{table}[H]
\centering
\caption{Tổng hợp 19 công trình tiêu biểu về dự báo bỏ học theo bộ tiêu chí thống nhất}
\label{tab:khaosat}
\renewcommand{\arraystretch}{1.25}
\footnotesize
\begin{tabularx}{\textwidth}{@{}p{2.4cm} p{1.6cm} p{2.6cm} X p{1.7cm}@{}}
\toprule
\textbf{Công trình} & \textbf{Nhãn} & \textbf{Phương pháp} & \textbf{Dữ liệu / Hiệu năng} & \textbf{Cửa sổ} \\
\midrule
Kloft 2014 \cite{Kloft2014} & Nhị phân & SVM, đặc trưng thủ công & MOOC tự thu thập; AUC theo tuần & Nhiều tuần \\
Taylor 2014 \cite{Taylor2014} & Nhị phân & Logistic Regression & MITx 6.002x & Tuần đầu \\
Fei \& Yeung 2015 \cite{Fei2015} & Nhị phân & IOHMM, RNN, LSTM & Coursera + edX; LSTM tốt nhất & Chuỗi tuần \\
Whitehill 2017 \cite{Whitehill2017} & Nhị phân & LR + đặc trưng thời gian & HarvardX, MITx; AUC $\approx$ 0,90 & Nhiều tuần \\
Nagrecha 2017 \cite{Nagrecha2017} & Nhị phân & Pipeline diễn giải & MITx, HarvardX & 1 cửa sổ \\
Vitiello 2018 \cite{Vitiello2018} & Nhị phân & LR, RF, AdaBoost & MOOC tự thu thập & Nhiều tuần \\
Alamri 2019 \cite{Alamri2019} & Nhị phân & RF, AdaBoost, GB, XGBoost & FutureLearn; Acc. 82--94\% & Tuần 1 \\
Xing \& Du 2019 \cite{Xing2019} & Nhị phân & Deep MLP cá nhân hóa & 2 nguồn MOOC & 1 cửa sổ \\
Mubarak 2021 \cite{Mubarak2021} & Nhị phân & CNN + LSTM (ma trận 30 ngày) & KDD Cup 2015; AUC $>$ 0,90 & Tuần đầu \\
Fu 2021 \cite{Fu2021} & Nhị phân & CLSA (CNN+LSTM+Att.) & KDD Cup 2015; AUC cao & 1 cửa sổ \\
Dass 2021 \cite{Dass2021} & Nhị phân & RF + tiến độ + SHAP & Open edX; Acc. 87,5\% & Theo ngày \\
Psathas 2021 \cite{Psathas2021} & Nhị phân & Nhiều bộ PL + tối ưu siêu TS & MOOC smart-city; Acc. $\approx$ 88\% & Sớm \\
Basnet 2022 \cite{Basnet2022} & Nhị phân & Đối sánh ML vs. DL & XuetangX, KDD; Acc. 84--88\% & 1 cửa sổ \\
Yu 2022 \cite{Yu2022} & Nhị phân & BiLSTM + meta-transfer & Nhiều MOOC; trans-domain & 1 cửa sổ \\
Martins 2023 \cite{Martins2023} & 3 mức & RF + xử lý mất cân bằng & ĐH Bồ Đào Nha & Phase đầu \\
Junejo 2025 \cite{Junejo2025} & 4 mức & 1D-CNN & OULAD; Acc. $>$ 93\% tại Q4 & Q1--Q4 \\
Aouarib 2025 \cite{Aouarib2025} & Nhị phân & GA + XGBoost & OULAD; Acc. 81--98\% & 4 cửa sổ \\
Benabbes 2025 \cite{Benabbes2025} & Nhị phân & 9 mô hình + SHAP, LIME & 17.883 SV; AUC 0,993 & 1 cửa sổ \\
Nguyen 2025 \cite{Nguyen2025} & 5 mức & GCN adversarial + GRU & MOOCCubeX; $\FoneD = 0{,}92$ & 4 tuần \\
\bottomrule
\end{tabularx}
\end{table}

Từ Bảng~\ref{tab:khaosat} có thể rút ra ba khoảng trống lớn, xác định không gian đóng góp cho đồ án:

\textit{Khoảng trống thứ nhất -- ngữ nghĩa rủi ro nhiều mức.} Khoảng 16 trong 19 công trình dùng nhãn nhị phân, và chỉ duy nhất một công trình (Nguyen và cộng sự) dùng nhãn 5 mức A/B/C/D/E trên MOOCCubeX. Nhãn nhị phân gộp chung nhóm ``đang hoạt động nhưng rủi ro cao'' với nhóm ``đã rời khóa'', mặc dù hai nhóm này có ý nghĩa can thiệp hoàn toàn khác nhau.

\textit{Khoảng trống thứ hai -- phân tích đa cửa sổ.} Chỉ khoảng 4 trong 19 công trình đánh giá tại nhiều cửa sổ quan sát; phần lớn còn lại cố định một thời điểm duy nhất. Điều này dẫn đến thiếu bằng chứng định lượng cho đánh đổi giữa thời điểm cảnh báo sớm và chất lượng dự báo, một vấn đề đã được Schwerter và cộng sự đặc biệt nhấn mạnh \cite{Schwerter2026}.

\textit{Khoảng trống thứ ba -- thiếu khung đối sánh thống nhất.} Trong toàn bộ 19 công trình, các trục thiết kế phụ trợ (điền giá trị thiếu, bộ trích xuất chuỗi, hàm mất mát) thường được cố định khi báo cáo kết quả, ít khi được benchmark song song có kiểm soát. Cả Dalipi và cộng sự lẫn Prenkaj và cộng sự đều nhấn mạnh khoảng trống này \cite{Dalipi2018,Prenkaj2020}.

Tổng hợp cả ba khoảng trống nêu trên chính là bối cảnh trực tiếp cho thiết kế của khung EarlyRisk-MOOC được trình bày trong Chương~\ref{ch:phuongphap}.

\section{Định vị thiết kế của khung đề xuất}
\label{sec:dinhvi}

Để làm rõ đóng góp về mặt thiết kế, Bảng~\ref{tab:dinhvi} định vị EarlyRisk-MOOC so với một số công trình tiêu biểu theo các tiêu chí liên quan đến cảnh báo sớm. Đây là so sánh định tính ở khía cạnh thiết kế; phần benchmark định lượng sẽ được trình bày trong Chương~\ref{ch:ketqua}.

\begin{table}[H]
\centering
\caption{Định vị thiết kế của EarlyRisk-MOOC so với các công trình khảo sát ($\checkmark$: có; $\triangle$: một phần; $\times$: không)}
\label{tab:dinhvi}
\renewcommand{\arraystretch}{1.3}
\footnotesize
\begin{tabularx}{\textwidth}{@{}p{3.1cm} c c c c c@{}}
\toprule
\textbf{Công trình} & \textbf{Cảnh báo sớm} & \textbf{Nhãn đa mức} & \textbf{Trục TK là biến TN} & \textbf{Đa cửa sổ} & \textbf{Báo cáo chi phí} \\
\midrule
Fei \& Yeung 2015 & $\checkmark$ & $\times$ & $\times$ & $\checkmark$ & $\times$ \\
Alamri 2019 & $\checkmark$ & $\times$ & $\times$ & $\times$ & $\triangle$ \\
Xing \& Du 2019 & $\checkmark$ & $\times$ & $\times$ & $\triangle$ & $\times$ \\
Mubarak 2021 & $\checkmark$ & $\times$ & $\triangle$ & $\times$ & $\times$ \\
Dass 2021 & $\checkmark$ & $\times$ & $\times$ & $\checkmark$ & $\triangle$ \\
Basnet 2022 & $\triangle$ & $\times$ & $\times$ & $\times$ & $\times$ \\
Junejo 2025 & $\checkmark$ & $\checkmark$ (4 mức) & $\times$ & $\checkmark$ & $\triangle$ \\
Aouarib 2025 & $\checkmark$ & $\times$ & $\times$ & $\checkmark$ & $\times$ \\
Nguyen 2025 & $\checkmark$ & $\checkmark$ (5 mức) & $\checkmark$ & $\times$ & $\checkmark$ \\
\textbf{EarlyRisk-MOOC} & $\checkmark$ & $\checkmark$ (5 mức) & $\checkmark$ & $\checkmark$ & $\checkmark$ \\
\bottomrule
\end{tabularx}
\end{table}

Như Bảng~\ref{tab:dinhvi} thể hiện, EarlyRisk-MOOC là khung duy nhất trong số các công trình khảo sát đáp ứng đồng thời cả năm tiêu chí: cảnh báo sớm trên nhiều cửa sổ, nhãn đa mức, coi các trục thiết kế là biến thí nghiệm, phân tích đa cửa sổ và báo cáo chi phí triển khai. Sự ``cân bằng'' này -- thay vì tối ưu một khía cạnh đơn lẻ -- chính là đặc trưng thiết kế cốt lõi của khung đề xuất.

\section{Kết luận chương}

Chương này đã khảo sát có hệ thống 19 công trình tiêu biểu thuộc hai họ phương pháp chính, phân tích hiện trạng và tổng hợp ba khoảng trống nghiên cứu: thiếu ngữ nghĩa rủi ro nhiều mức, thiếu phân tích đa cửa sổ, và thiếu khung đối sánh thống nhất. Ba khoảng trống này định hướng trực tiếp cho thiết kế của khung EarlyRisk-MOOC. Chương tiếp theo trình bày cơ sở lý thuyết và kiến trúc chi tiết của khung đề xuất.

\clearpage
